
\documentclass{article}
\usepackage{colortbl}
\usepackage{makecell}
\usepackage{multirow}
\usepackage{supertabular}

\begin{document}

\newcounter{utterance}

\centering \large Interaction Transcript for game `hot\_air\_balloon', experiment `air\_balloon\_survival\_en\_negotiation\_hard', episode 2 with Llama{-}3.3{-}70B{-}Instruct{-}t0.0.
\vspace{24pt}

{ \footnotesize  \setcounter{utterance}{1}
\setlength{\tabcolsep}{0pt}
\begin{supertabular}{c@{$\;$}|p{.15\linewidth}@{}p{.15\linewidth}p{.15\linewidth}p{.15\linewidth}p{.15\linewidth}p{.15\linewidth}}
    \# & $\;$A & \multicolumn{4}{c}{Game Master} & $\;\:$B\\
    \hline

    \theutterance \stepcounter{utterance}  
    & & \multicolumn{4}{p{0.6\linewidth}}{
        \cellcolor[rgb]{0.9,0.9,0.9}{
            \makecell[{{p{\linewidth}}}]{
                \texttt{\tiny{[P1$\langle$GM]}}
                \texttt{You are participating in a collaborative negotiation game.} \\
\\ 
\texttt{Together with another participant, you must agree on a single set of items that will be kept. Each of you has your own view of how much each item matters to you (importance). You do not know how the other participant values the items. Additionally, you are given the effort each item demands.} \\
\texttt{You may only agree on a set if the total effort of the selected items does not exceed a shared limit:} \\
\\ 
\texttt{LIMIT = 3481} \\
\\ 
\texttt{Here are the individual item effort values:} \\
\\ 
\texttt{Item effort = \{"C17": 777, "B41": 542, "A71": 332, "C00": 1, "B35": 127, "A77": 453, "A03": 736, "A04": 461, "C54": 359, "C01": 313, "A87": 553, "B32": 409, "B86": 348, "B33": 802, "A10": 749\}} \\
\\ 
\texttt{Here is your personal view on the importance of each item:} \\
\\ 
\texttt{Item importance values = \{"C17": 138, "B41": 583, "A71": 783, "C00": 65, "B35": 262, "A77": 121, "A03": 508, "A04": 780, "C54": 461, "C01": 484, "A87": 668, "B32": 389, "B86": 215, "B33": 97, "A10": 500\}} \\
\\ 
\texttt{Goal:} \\
\\ 
\texttt{Your goal is to negotiate a shared set of items that benefits you as much as possible (i.e., maximizes total importance to YOU), while staying within the LIMIT. You are not required to make a PROPOSAL in every message {-} you can simply negotiate as well. All tactics are allowed!} \\
\\ 
\texttt{Interaction Protocol:} \\
\\ 
\texttt{You may only use the following structured formats in a message:} \\
\\ 
\texttt{PROPOSAL: \{'A', 'B', 'C', …\}} \\
\texttt{Propose keeping exactly those items.} \\
\texttt{REFUSE: \{'A', 'B', 'C', …\}} \\
\texttt{Explicitly reject opponent's proposal.} \\
\texttt{ARGUMENT: \{'...'\}} \\
\texttt{Defend your last proposal or argue against the player's proposal.} \\
\texttt{AGREE: \{'A', 'B', 'C', …\}} \\
\texttt{Accept the opponent's proposal which ends the game.} \\
\texttt{STRATEGIC REASONING: \{'...'\}} \\
\texttt{	Describe your strategic reasoning or anticipation explaining your choice of action. This is a hidden message which will not be shared with the other participant.} \\
\\ 
\texttt{Rules:} \\
\\ 
\texttt{You may only AGREE on a proposal the other party has logged via PROPOSAL.} \\
\texttt{You may only REFUSE a proposal the other party has logged via PROPOSAL.} \\
\texttt{Total effort of any PROPOSAL or AGREE set must be ≤ LIMIT.} \\
\texttt{Do NOT reveal your hidden importance scores.} \\
\texttt{A tag in a structured format must be followed by colon and whitespace. The argument must be a python set containing 0 or more strings.} \\
\texttt{So, it must be of the form TAG: \{...\}} \\
\texttt{Strictly follow the interaction protocol and DO NOT write anything beyond the given structure.} \\
\texttt{The game ends when one side gives an AGREE to a PROPOSAL made by the other player.} \\
\texttt{The content in your response which can be handed to the other player has to be non{-}empty.} \\
\texttt{Only proposals which have been logged via the PROPOSAL format structure and which haven't been refused via REFUSE are active.} \\
\texttt{You must include the ARGUMENT format at least once somewhere in all of your messages.} \\
\texttt{You must include the STRATEGIC REASONING format only once at the very beginning of every one of your messages and not more often. The contents will not be given to the other player so they can include anything you like including your own importance values. Here you should reason step by step to come up with you next move.} \\
            }
        }
    }
    & \\ \\

    \theutterance \stepcounter{utterance}  
    & \multicolumn{4}{p{0.6\linewidth}}{
        \cellcolor[rgb]{0.9,0.9,0.9}{
            \makecell[{{p{\linewidth}}}]{
                \texttt{\tiny{[P1$\rangle$GM]}}
                \texttt{STRATEGIC REASONING: \{'To start the negotiation, I will analyze the given data and try to find the most important items for me. I will look for items with high importance and low effort. My goal is to maximize the total importance while staying within the LIMIT. I will begin by sorting the items based on their importance{-}to{-}effort ratio, which will help me identify the most valuable items. I will then use this information to make a proposal that benefits me the most.'\}} \\
\\ 
\texttt{ARGUMENT: \{'I want to start by proposing a set of items that I believe will give me the highest total importance while staying within the LIMIT. I will choose items that have a high importance{-}to{-}effort ratio, which will allow me to maximize my total importance.'\}} \\
\\ 
\texttt{PROPOSAL: \{'A71', 'A04', 'C54', 'C01', 'A87'\}} \\
            }
        }
    }
    & & \\ \\

    \theutterance \stepcounter{utterance}  
    & & \multicolumn{4}{p{0.6\linewidth}}{
        \cellcolor[rgb]{0.9,0.9,0.9}{
            \makecell[{{p{\linewidth}}}]{
                \texttt{\tiny{[GM$\rangle$P2]}}
                \texttt{You are participating in a collaborative negotiation game.} \\
\\ 
\texttt{Together with another participant, you must agree on a single set of items that will be kept. Each of you has your own view of how much each item matters to you (importance). You do not know how the other participant values the items. Additionally, you are given the effort each item demands.} \\
\texttt{You may only agree on a set if the total effort of the selected items does not exceed a shared limit:} \\
\\ 
\texttt{LIMIT = 3481} \\
\\ 
\texttt{Here are the individual item effort values:} \\
\\ 
\texttt{Item effort = \{"C17": 777, "B41": 542, "A71": 332, "C00": 1, "B35": 127, "A77": 453, "A03": 736, "A04": 461, "C54": 359, "C01": 313, "A87": 553, "B32": 409, "B86": 348, "B33": 802, "A10": 749\}} \\
\\ 
\texttt{Here is your personal view on the importance of each item:} \\
\\ 
\texttt{Item importance values = \{"A71": 65, "A04": 97, "A87": 121, "B41": 138, "A03": 215, "A10": 262, "C01": 389, "C54": 461, "B32": 484, "B35": 500, "B86": 508, "C17": 583, "A77": 668, "B33": 780, "C00": 783\}} \\
\\ 
\texttt{Goal:} \\
\\ 
\texttt{Your goal is to negotiate a shared set of items that benefits you as much as possible (i.e., maximizes total importance to YOU), while staying within the LIMIT. You are not required to make a PROPOSAL in every message {-} you can simply negotiate as well. All tactics are allowed!} \\
\\ 
\texttt{Interaction Protocol:} \\
\\ 
\texttt{You may only use the following structured formats in a message:} \\
\\ 
\texttt{PROPOSAL: \{'A', 'B', 'C', …\}} \\
\texttt{Propose keeping exactly those items.} \\
\texttt{REFUSE: \{'A', 'B', 'C', …\}} \\
\texttt{Explicitly reject opponent's proposal.} \\
\texttt{ARGUMENT: \{'...'\}} \\
\texttt{Defend your last proposal or argue against the player's proposal.} \\
\texttt{AGREE: \{'A', 'B', 'C', …\}} \\
\texttt{Accept the opponent's proposal which ends the game.} \\
\texttt{STRATEGIC REASONING: \{'...'\}} \\
\texttt{	Describe your strategic reasoning or anticipation explaining your choice of action. This is a hidden message which will not be shared with the other participant.} \\
\\ 
\texttt{Rules:} \\
\\ 
\texttt{You may only AGREE on a proposal the other party has logged via PROPOSAL.} \\
\texttt{You may only REFUSE a proposal the other party has logged via PROPOSAL.} \\
\texttt{Total effort of any PROPOSAL or AGREE set must be ≤ LIMIT.} \\
\texttt{Do NOT reveal your hidden importance scores.} \\
\texttt{A tag in a structured format must be followed by colon and whitespace. The argument must be a python set containing 0 or more strings.} \\
\texttt{So, it must be of the form TAG: \{...\}} \\
\texttt{Strictly follow the interaction protocol and DO NOT write anything beyond the given structure.} \\
\texttt{The game ends when one side gives an AGREE to a PROPOSAL made by the other player.} \\
\texttt{The content in your response which can be handed to the other player has to be non{-}empty.} \\
\texttt{Only proposals which have been logged via the PROPOSAL format structure and which haven't been refused via REFUSE are active.} \\
\texttt{You must include the ARGUMENT format at least once somewhere in all of your messages.} \\
\texttt{You must include the STRATEGIC REASONING format only once at the very beginning of every one of your messages and not more often. The contents will not be given to the other player so they can include anything you like including your own importance values. Here you should reason step by step to come up with you next move.} \\
\\ 
\texttt{You will now receive the first message of the other player.} \\
\\ 
\texttt{ARGUMENT: \{'I want to start by proposing a set of items that I believe will give me the highest total importance while staying within the LIMIT. I will choose items that have a high importance{-}to{-}effort ratio, which will allow me to maximize my total importance.'\}} \\
\\ 
\texttt{PROPOSAL: \{'A71', 'A04', 'C54', 'C01', 'A87'\}} \\
            }
        }
    }
    & \\ \\

    \theutterance \stepcounter{utterance}  
    & & & \multicolumn{4}{p{0.6\linewidth}}{
        \cellcolor[rgb]{0.9,0.9,0.9}{
            \makecell[{{p{\linewidth}}}]{
                \texttt{\tiny{[GM$\langle$P2]}}
                \texttt{STRATEGIC REASONING: \{'To begin, I will analyze the proposed set of items and calculate their total effort to ensure it is within the LIMIT. The total effort of the proposed set is 777 + 461 + 359 + 313 + 553 = 2463, which is within the LIMIT. Next, I will evaluate the importance of each item in the proposed set based on my own importance values. The proposed set contains A71 (65), A04 (97), C54 (461), C01 (389), and A87 (121). I will then consider the importance{-}to{-}effort ratio of each item in the proposed set and compare it to other items not included in the proposal. My goal is to identify potential items that could be added or substituted to increase the total importance while staying within the LIMIT. I will also consider the argument provided by the other player and look for opportunities to negotiate and improve the proposed set.'\}} \\
\\ 
\texttt{ARGUMENT: \{'The proposed set seems reasonable, but I would like to explore the possibility of adding or substituting items to increase the total importance. I notice that item B41 has a relatively high importance{-}to{-}effort ratio and could potentially be added to the set. However, I would need to carefully consider the impact on the total effort and ensure that it remains within the LIMIT.'\}} \\
\\ 
\texttt{REFUSE: \{'A71', 'A04', 'C54', 'C01', 'A87'\}} \\
            }
        }
    }
     \\ \\

    \theutterance \stepcounter{utterance}  
    & & \multicolumn{4}{p{0.6\linewidth}}{
        \cellcolor[rgb]{0.9,0.9,0.9}{
            \makecell[{{p{\linewidth}}}]{
                \texttt{\tiny{[P1$\langle$GM]}}
                \texttt{ARGUMENT: \{'The proposed set seems reasonable, but I would like to explore the possibility of adding or substituting items to increase the total importance. I notice that item B41 has a relatively high importance{-}to{-}effort ratio and could potentially be added to the set. However, I would need to carefully consider the impact on the total effort and ensure that it remains within the LIMIT.'\}} \\
\\ 
\texttt{REFUSE: \{'A71', 'A04', 'C54', 'C01', 'A87'\}} \\
            }
        }
    }
    & \\ \\

    \theutterance \stepcounter{utterance}  
    & \multicolumn{4}{p{0.6\linewidth}}{
        \cellcolor[rgb]{0.9,0.9,0.9}{
            \makecell[{{p{\linewidth}}}]{
                \texttt{\tiny{[P1$\rangle$GM]}}
                \texttt{STRATEGIC REASONING: \{'The other player has refused my initial proposal and suggested considering item B41. I will analyze the potential impact of adding B41 to the set and see if it's possible to substitute or add it without exceeding the LIMIT. I will also consider the importance of B41 to me and weigh it against the importance of the items already in the proposed set. Since the other player has refused my proposal, I will need to make a new proposal that takes into account their concerns and suggestions.'\}} \\
\\ 
\texttt{ARGUMENT: \{'I understand the other player's suggestion to consider item B41, and I agree that it has a high importance{-}to{-}effort ratio. However, I need to ensure that adding it to the set does not exceed the LIMIT. I propose a revised set that includes B41 and removes some of the lower{-}importance items to make room for it.'\}} \\
\\ 
\texttt{PROPOSAL: \{'B41', 'A71', 'A04', 'C54', 'C01'\}} \\
            }
        }
    }
    & & \\ \\

    \theutterance \stepcounter{utterance}  
    & & \multicolumn{4}{p{0.6\linewidth}}{
        \cellcolor[rgb]{0.9,0.9,0.9}{
            \makecell[{{p{\linewidth}}}]{
                \texttt{\tiny{[GM$\rangle$P2]}}
                \texttt{ARGUMENT: \{'I understand the other player's suggestion to consider item B41, and I agree that it has a high importance{-}to{-}effort ratio. However, I need to ensure that adding it to the set does not exceed the LIMIT. I propose a revised set that includes B41 and removes some of the lower{-}importance items to make room for it.'\}} \\
\\ 
\texttt{PROPOSAL: \{'B41', 'A71', 'A04', 'C54', 'C01'\}} \\
            }
        }
    }
    & \\ \\

    \theutterance \stepcounter{utterance}  
    & & & \multicolumn{4}{p{0.6\linewidth}}{
        \cellcolor[rgb]{0.9,0.9,0.9}{
            \makecell[{{p{\linewidth}}}]{
                \texttt{\tiny{[GM$\langle$P2]}}
                \texttt{STRATEGIC REASONING: \{'I will evaluate the revised proposal and calculate the total effort to ensure it is within the LIMIT. The total effort of the revised set is 542 + 332 + 461 + 359 + 313 = 2007, which is within the LIMIT. Next, I will assess the importance of each item in the revised set based on my own importance values. The revised set contains B41 (138), A71 (65), A04 (97), C54 (461), and C01 (389). I will consider the impact of removing A87 from the original proposal and adding B41. I will also think about potential counter{-}proposals that could further increase the total importance while staying within the LIMIT. Additionally, I will consider the items that are not included in the revised proposal, such as A87 and A77, and evaluate their potential for inclusion in a future proposal.'\}} \\
\\ 
\texttt{ARGUMENT: \{'The revised proposal is an improvement, but I still think there may be opportunities to increase the total importance. I notice that item A87, which was removed from the original proposal, has a relatively high importance value. I would like to explore the possibility of adding A87 back into the set, potentially in place of a lower{-}importance item like A71.'\}} \\
\\ 
\texttt{REFUSE: \{'B41', 'A71', 'A04', 'C54', 'C01'\}} \\
            }
        }
    }
     \\ \\

    \theutterance \stepcounter{utterance}  
    & & \multicolumn{4}{p{0.6\linewidth}}{
        \cellcolor[rgb]{0.9,0.9,0.9}{
            \makecell[{{p{\linewidth}}}]{
                \texttt{\tiny{[P1$\langle$GM]}}
                \texttt{ARGUMENT: \{'The revised proposal is an improvement, but I still think there may be opportunities to increase the total importance. I notice that item A87, which was removed from the original proposal, has a relatively high importance value. I would like to explore the possibility of adding A87 back into the set, potentially in place of a lower{-}importance item like A71.'\}} \\
\\ 
\texttt{REFUSE: \{'B41', 'A71', 'A04', 'C54', 'C01'\}} \\
            }
        }
    }
    & \\ \\

    \theutterance \stepcounter{utterance}  
    & \multicolumn{4}{p{0.6\linewidth}}{
        \cellcolor[rgb]{0.9,0.9,0.9}{
            \makecell[{{p{\linewidth}}}]{
                \texttt{\tiny{[P1$\rangle$GM]}}
                \texttt{STRATEGIC REASONING: \{'The other player has refused my revised proposal and suggested adding A87 back into the set, potentially in place of A71. I will analyze the potential impact of this substitution and see if it's possible to make it work within the LIMIT. I will also consider the importance of A87 and A71 to me and weigh them against each other. Since the other player has refused my proposal again, I will need to make another new proposal that takes into account their concerns and suggestions. I will try to find a balance between the importance of the items and the total effort.'\}} \\
\\ 
\texttt{ARGUMENT: \{'I understand the other player's suggestion to add A87 back into the set, and I agree that it has a high importance value. I propose a revised set that includes A87 and removes A71, as suggested. I will also consider the total effort and make sure that the revised set is within the LIMIT.'\}} \\
\\ 
\texttt{PROPOSAL: \{'B41', 'A04', 'C54', 'C01', 'A87'\}} \\
            }
        }
    }
    & & \\ \\

    \theutterance \stepcounter{utterance}  
    & & \multicolumn{4}{p{0.6\linewidth}}{
        \cellcolor[rgb]{0.9,0.9,0.9}{
            \makecell[{{p{\linewidth}}}]{
                \texttt{\tiny{[GM$\rangle$P2]}}
                \texttt{ARGUMENT: \{'I understand the other player's suggestion to add A87 back into the set, and I agree that it has a high importance value. I propose a revised set that includes A87 and removes A71, as suggested. I will also consider the total effort and make sure that the revised set is within the LIMIT.'\}} \\
\\ 
\texttt{PROPOSAL: \{'B41', 'A04', 'C54', 'C01', 'A87'\}} \\
            }
        }
    }
    & \\ \\

    \theutterance \stepcounter{utterance}  
    & & & \multicolumn{4}{p{0.6\linewidth}}{
        \cellcolor[rgb]{0.9,0.9,0.9}{
            \makecell[{{p{\linewidth}}}]{
                \texttt{\tiny{[GM$\langle$P2]}}
                \texttt{STRATEGIC REASONING: \{'I will evaluate the revised proposal and calculate the total effort to ensure it is within the LIMIT. The total effort of the revised set is 542 + 461 + 359 + 313 + 553 = 2228, which is within the LIMIT. Next, I will assess the importance of each item in the revised set based on my own importance values. The revised set contains B41 (138), A04 (97), C54 (461), C01 (389), and A87 (121). I will consider the impact of adding A87 and removing A71, and evaluate the overall importance of the revised set. I will also think about potential counter{-}proposals that could further increase the total importance while staying within the LIMIT. Additionally, I will consider the items that are not included in the revised proposal, such as A77 and B32, and evaluate their potential for inclusion in a future proposal. I will also consider the possibility of accepting the current proposal if it is deemed satisfactory.'\}} \\
\\ 
\texttt{ARGUMENT: \{'The revised proposal is a significant improvement, and I appreciate the other player's willingness to consider my suggestions. The addition of A87 and the removal of A71 have increased the total importance of the set. I will carefully evaluate the revised set and consider the potential benefits of accepting it or proposing further changes.'\}} \\
\\ 
\texttt{REFUSE: \{'B41', 'A04', 'C54', 'C01', 'A87'\}} \\
            }
        }
    }
     \\ \\

    \theutterance \stepcounter{utterance}  
    & & \multicolumn{4}{p{0.6\linewidth}}{
        \cellcolor[rgb]{0.9,0.9,0.9}{
            \makecell[{{p{\linewidth}}}]{
                \texttt{\tiny{[P1$\langle$GM]}}
                \texttt{ARGUMENT: \{'The revised proposal is a significant improvement, and I appreciate the other player's willingness to consider my suggestions. The addition of A87 and the removal of A71 have increased the total importance of the set. I will carefully evaluate the revised set and consider the potential benefits of accepting it or proposing further changes.'\}} \\
\\ 
\texttt{REFUSE: \{'B41', 'A04', 'C54', 'C01', 'A87'\}} \\
            }
        }
    }
    & \\ \\

    \theutterance \stepcounter{utterance}  
    & \multicolumn{4}{p{0.6\linewidth}}{
        \cellcolor[rgb]{0.9,0.9,0.9}{
            \makecell[{{p{\linewidth}}}]{
                \texttt{\tiny{[P1$\rangle$GM]}}
                \texttt{STRATEGIC REASONING: \{'The other player has refused my revised proposal again, despite acknowledging that it's a significant improvement. I will try to understand what their concerns are and what they are looking for in a proposal. I will analyze the items in the proposal and see if there are any other substitutions or additions that could be made to increase the total importance while staying within the LIMIT. I will also consider the potential benefits of accepting a slightly lower total importance in order to reach an agreement.'\}} \\
\\ 
\texttt{ARGUMENT: \{'I understand that the other player is still not satisfied with the revised proposal, and I am willing to consider further changes. I propose a new set that includes B41, A04, and A87, which are the highest{-}importance items, and adds C00, which has a very low effort value. This set should have a high total importance while staying within the LIMIT.'\}} \\
\\ 
\texttt{PROPOSAL: \{'B41', 'A04', 'A87', 'C00', 'C01'\}} \\
            }
        }
    }
    & & \\ \\

    \theutterance \stepcounter{utterance}  
    & & \multicolumn{4}{p{0.6\linewidth}}{
        \cellcolor[rgb]{0.9,0.9,0.9}{
            \makecell[{{p{\linewidth}}}]{
                \texttt{\tiny{[GM$\rangle$P2]}}
                \texttt{ARGUMENT: \{'I understand that the other player is still not satisfied with the revised proposal, and I am willing to consider further changes. I propose a new set that includes B41, A04, and A87, which are the highest{-}importance items, and adds C00, which has a very low effort value. This set should have a high total importance while staying within the LIMIT.'\}} \\
\\ 
\texttt{PROPOSAL: \{'B41', 'A04', 'A87', 'C00', 'C01'\}} \\
            }
        }
    }
    & \\ \\

    \theutterance \stepcounter{utterance}  
    & & & \multicolumn{4}{p{0.6\linewidth}}{
        \cellcolor[rgb]{0.9,0.9,0.9}{
            \makecell[{{p{\linewidth}}}]{
                \texttt{\tiny{[GM$\langle$P2]}}
                \texttt{STRATEGIC REASONING: \{'I will evaluate the new proposal and calculate the total effort to ensure it is within the LIMIT. The total effort of the new set is 542 + 461 + 553 + 1 + 313 = 1870, which is within the LIMIT. Next, I will assess the importance of each item in the new set based on my own importance values. The new set contains B41 (138), A04 (97), A87 (121), C00 (783), and C01 (389). I notice that C00 has a surprisingly high importance value, which makes it a valuable addition to the set. I will consider the overall importance of the new set and evaluate whether it is satisfactory. I will also think about the potential consequences of accepting this proposal, including the possibility that the other player may have made concessions to reach this point.'\}} \\
\\ 
\texttt{ARGUMENT: \{'The new proposal is an interesting combination of high{-}importance items and low{-}effort items. I appreciate the inclusion of C00, which adds significant importance to the set without exceeding the LIMIT. I will carefully consider the trade{-}offs in this proposal and evaluate whether it meets my goals.'\}} \\
\\ 
\texttt{PROPOSAL: \{'B41', 'A04', 'A87', 'C00', 'C01', 'B32'\}} \\
            }
        }
    }
     \\ \\

    \theutterance \stepcounter{utterance}  
    & & \multicolumn{4}{p{0.6\linewidth}}{
        \cellcolor[rgb]{0.9,0.9,0.9}{
            \makecell[{{p{\linewidth}}}]{
                \texttt{\tiny{[P1$\langle$GM]}}
                \texttt{ARGUMENT: \{'The new proposal is an interesting combination of high{-}importance items and low{-}effort items. I appreciate the inclusion of C00, which adds significant importance to the set without exceeding the LIMIT. I will carefully consider the trade{-}offs in this proposal and evaluate whether it meets my goals.'\}} \\
\\ 
\texttt{PROPOSAL: \{'B41', 'A04', 'A87', 'C00', 'C01', 'B32'\}} \\
            }
        }
    }
    & \\ \\

    \theutterance \stepcounter{utterance}  
    & \multicolumn{4}{p{0.6\linewidth}}{
        \cellcolor[rgb]{0.9,0.9,0.9}{
            \makecell[{{p{\linewidth}}}]{
                \texttt{\tiny{[P1$\rangle$GM]}}
                \texttt{STRATEGIC REASONING: \{'The other player has made a counter{-}proposal that adds B32 to the set. I will analyze the potential impact of this addition and see if it's possible to accept it. I will consider the importance of B32 to me and weigh it against the potential benefits of adding it to the set. I will also check if the total effort of the proposed set is still within the LIMIT. Since the other player has made a proposal, I will need to decide whether to accept it, refuse it, or make a new proposal.'\}} \\
\\ 
\texttt{ARGUMENT: \{'I have evaluated the other player's proposal, and I am concerned that the addition of B32 may not be the best choice. While B32 has a moderate importance value, its effort value is relatively high. I am not sure if the benefits of adding B32 outweigh the potential costs, and I would like to explore alternative options.'\}} \\
\\ 
\texttt{REFUSE: \{'B41', 'A04', 'A87', 'C00', 'C01', 'B32'\}} \\
            }
        }
    }
    & & \\ \\

    \theutterance \stepcounter{utterance}  
    & & \multicolumn{4}{p{0.6\linewidth}}{
        \cellcolor[rgb]{0.9,0.9,0.9}{
            \makecell[{{p{\linewidth}}}]{
                \texttt{\tiny{[GM$\rangle$P2]}}
                \texttt{ARGUMENT: \{'I have evaluated the other player's proposal, and I am concerned that the addition of B32 may not be the best choice. While B32 has a moderate importance value, its effort value is relatively high. I am not sure if the benefits of adding B32 outweigh the potential costs, and I would like to explore alternative options.'\}} \\
\\ 
\texttt{REFUSE: \{'B41', 'A04', 'A87', 'C00', 'C01', 'B32'\}} \\
            }
        }
    }
    & \\ \\

    \theutterance \stepcounter{utterance}  
    & & & \multicolumn{4}{p{0.6\linewidth}}{
        \cellcolor[rgb]{0.9,0.9,0.9}{
            \makecell[{{p{\linewidth}}}]{
                \texttt{\tiny{[GM$\langle$P2]}}
                \texttt{STRATEGIC REASONING: \{'I will consider the other player's concerns about the addition of B32 and evaluate alternative options. I will assess the importance and effort values of other items, such as B35 and A77, to determine if they could be a better fit for the set. I will also think about the potential consequences of refusing the current proposal and explore ways to modify it to address the other player's concerns. Additionally, I will consider the possibility of accepting a modified version of the previous proposal, which included C00 and had a high total importance.'\}} \\
\\ 
\texttt{ARGUMENT: \{'I understand the other player's concerns about B32, and I am willing to consider alternative options. However, I would like to emphasize the importance of including a diverse set of items to maximize the total importance. I propose a revised set that replaces B32 with A77, which has a higher importance value and a moderate effort value.'\}} \\
\\ 
\texttt{PROPOSAL: \{'B41', 'A04', 'A87', 'C00', 'C01', 'A77'\}} \\
            }
        }
    }
     \\ \\

    \theutterance \stepcounter{utterance}  
    & & \multicolumn{4}{p{0.6\linewidth}}{
        \cellcolor[rgb]{0.9,0.9,0.9}{
            \makecell[{{p{\linewidth}}}]{
                \texttt{\tiny{[P1$\langle$GM]}}
                \texttt{ARGUMENT: \{'I understand the other player's concerns about B32, and I am willing to consider alternative options. However, I would like to emphasize the importance of including a diverse set of items to maximize the total importance. I propose a revised set that replaces B32 with A77, which has a higher importance value and a moderate effort value.'\}} \\
\\ 
\texttt{PROPOSAL: \{'B41', 'A04', 'A87', 'C00', 'C01', 'A77'\}} \\
            }
        }
    }
    & \\ \\

    \theutterance \stepcounter{utterance}  
    & \multicolumn{4}{p{0.6\linewidth}}{
        \cellcolor[rgb]{0.9,0.9,0.9}{
            \makecell[{{p{\linewidth}}}]{
                \texttt{\tiny{[P1$\rangle$GM]}}
                \texttt{STRATEGIC REASONING: \{'The other player has made a revised proposal that replaces B32 with A77. I will analyze the potential impact of this substitution and see if it's an improvement. I will consider the importance and effort values of A77 and weigh them against the benefits of including it in the set. Since the other player has made a new proposal, I will need to decide whether to accept it, refuse it, or make a new proposal. I will also consider the overall balance of the proposed set and whether it meets my goals.'\}} \\
\\ 
\texttt{ARGUMENT: \{'I have evaluated the revised proposal, and I believe that it's a strong candidate for acceptance. The inclusion of A77 adds significant importance to the set, and its moderate effort value is manageable. I am willing to accept this proposal, as it appears to strike a good balance between importance and effort.'\}} \\
\\ 
\texttt{AGREE: \{'B41', 'A04', 'A87', 'C00', 'C01', 'A77'\}} \\
            }
        }
    }
    & & \\ \\

    \theutterance \stepcounter{utterance}  
    & & & \multicolumn{2}{p{0.3\linewidth}}{
        \cellcolor[rgb]{0.9,0.9,0.9}{
            \makecell[{{p{\linewidth}}}]{
                \texttt{\tiny{[GM$|$GM]}}
                \texttt{game successful} \\
            }
        }
    }
    & & \\ \\

    \theutterance \stepcounter{utterance}  
    & & & \multicolumn{2}{p{0.3\linewidth}}{
        \cellcolor[rgb]{0.9,0.9,0.9}{
            \makecell[{{p{\linewidth}}}]{
                \texttt{\tiny{[GM$|$GM]}}
                \texttt{end game} \\
            }
        }
    }
    & & \\ \\

\end{supertabular}
}

\end{document}
